% !TEX program = xelatex
%
% Epiphany --- Reference Suite (companion specification)
% Companion to the Core Specification. Compile with XeLaTeX.
%
% This document is versioned independently of the Core Specification
% (independent semver; see the Versioning note in the front matter). Its preamble
% is intentionally a self-contained copy of the core specification's preamble so
% the two documents build independently; factoring a shared preamble file is a
% later cleanup, not a v0.1 deliverable.

\documentclass[11pt,letterpaper]{report}

% ---------------------------------------------------------------------------
% Packages
% ---------------------------------------------------------------------------
\usepackage{fontspec}
\usepackage{geometry}
\geometry{
  letterpaper,
  top=1.05in,
  bottom=1.05in,
  left=1.15in,
  right=1.15in,
  headheight=15pt
}

\usepackage[english]{babel}
\usepackage{microtype}
\usepackage{parskip}
\usepackage{xcolor}
\usepackage{hyperref}
\usepackage{enumitem}
\usepackage{titlesec}
\usepackage{fancyhdr}
\usepackage{booktabs}
\usepackage{array}
\usepackage{longtable}
\usepackage{listings}
\usepackage{amsmath}
\usepackage{amssymb}
\usepackage{tcolorbox}
\tcbuselibrary{breakable, skins}

% ---------------------------------------------------------------------------
% Color palette (shared with the core specification)
% ---------------------------------------------------------------------------
\definecolor{epiphanyteal}{HTML}{1A4044}
\definecolor{epiphanygold}{HTML}{8E6E2E}
\definecolor{epiphanyink}{HTML}{1F1B16}
\definecolor{epiphanyslate}{HTML}{6B6660}
\definecolor{epiphanycream}{HTML}{F8F4ED}
\definecolor{epiphanymist}{HTML}{ECE8E0}
\definecolor{epiphanycode}{HTML}{2A2520}
\definecolor{epiphanycrimson}{HTML}{7A2424}

\hypersetup{
  colorlinks=true,
  linkcolor=epiphanyteal,
  citecolor=epiphanyteal,
  urlcolor=epiphanygold,
  pdftitle={Epiphany --- Reference Suite},
  pdfauthor={The Epiphany Project},
  pdfsubject={Reference Suite companion for the Epiphany music notation platform},
  pdfkeywords={music notation, engraving, reference suite, conformance, solver tiers, test scores},
  bookmarksnumbered=true,
  bookmarksopen=true
}

% ---------------------------------------------------------------------------
% Typography (shared with the core specification)
% ---------------------------------------------------------------------------
\setmainfont{TeX Gyre Pagella}[Numbers={OldStyle, Proportional}, Ligatures={TeX, Common}]
\setsansfont{TeX Gyre Heros}[Scale=0.94, Ligatures={TeX, Common}]
\setmonofont{TeX Gyre Cursor}[Scale=0.88, Ligatures={TeX}]
\newfontfamily\titlefont{TeX Gyre Pagella}[Numbers={OldStyle}, Ligatures={TeX, Common}]
\newcommand{\tablenums}[1]{{\addfontfeatures{Numbers={Lining,Tabular}}#1}}
\newcommand{\sectionsc}[1]{{\addfontfeatures{Letters=SmallCaps}#1}}

% ---------------------------------------------------------------------------
% Section styling (shared with the core specification)
% ---------------------------------------------------------------------------
\titleformat{\chapter}[display]
  {\normalfont\filright}
  {\raggedright\color{epiphanygold}\fontsize{14pt}{16pt}\selectfont
    \scshape Chapter\ \thechapter}
  {16pt}
  {\raggedright\color{epiphanyteal}\fontsize{32pt}{36pt}\selectfont\bfseries}
  [\vspace{4pt}{\color{epiphanygold}\rule{2in}{0.6pt}}]
\titlespacing*{\chapter}{0pt}{-20pt}{30pt}
\titleformat{\section}
  {\normalfont\Large\bfseries\color{epiphanyteal}}
  {\color{epiphanygold}\thesection}{1em}{}
\titleformat{\subsection}
  {\normalfont\large\bfseries\color{epiphanyteal}}
  {\color{epiphanygold}\thesubsection}{1em}{}
\titleformat{\subsubsection}
  {\normalfont\normalsize\bfseries\color{epiphanyink}}
  {\thesubsubsection}{1em}{}

% ---------------------------------------------------------------------------
% Headers and footers (shared with the core specification)
% ---------------------------------------------------------------------------
\pagestyle{fancy}
\fancyhf{}
\renewcommand{\headrulewidth}{0pt}
\renewcommand{\footrulewidth}{0pt}
\fancyhead[L]{\small\scshape\color{epiphanyslate}Epiphany --- Reference Suite}
\fancyhead[R]{\small\itshape\color{epiphanyslate}\leftmark}
\fancyfoot[C]{\small\color{epiphanyslate}\thepage}
\renewcommand{\headrule}{
  \color{epiphanygold!50}\hrule width\headwidth height 0.4pt
  \vspace{1pt}
  \color{epiphanygold!30}\hrule width\headwidth height 0.2pt
}

% ---------------------------------------------------------------------------
% Code listing style (shared with the core specification)
% ---------------------------------------------------------------------------
\lstdefinelanguage{Rust}{
  keywords={fn,let,mut,pub,struct,enum,impl,trait,for,in,if,else,match,return,
           use,mod,crate,self,Self,as,where,move,async,await,const,static,
           ref,type,unsafe,extern,dyn,box,break,continue,loop,while},
  keywordstyle=\color{epiphanyteal}\bfseries,
  ndkeywords={i8,i16,i32,i64,i128,u8,u16,u32,u64,u128,f32,f64,bool,char,str,
              String,Vec,Option,Result,Box,Rc,Arc,HashMap,BTreeMap,
              NonZeroU16,NonZeroU32,NonZeroU64,Duration,Timestamp},
  ndkeywordstyle=\color{epiphanygold}\bfseries,
  sensitive=true,
  comment=[l]{//},
  morecomment=[s]{/*}{*/},
  commentstyle=\color{epiphanyslate}\itshape,
  stringstyle=\color{epiphanycrimson},
  morestring=[b]",
  morestring=[b]'
}
\lstset{
  basicstyle=\ttfamily\small\color{epiphanycode},
  backgroundcolor=\color{epiphanycream},
  frame=leftline,
  rulecolor=\color{epiphanygold!60},
  framesep=8pt,
  framerule=1.5pt,
  xleftmargin=10pt,
  xrightmargin=4pt,
  breaklines=true,
  showstringspaces=false,
  numberstyle=\tiny\color{epiphanyslate},
  numbersep=10pt,
  captionpos=b,
  aboveskip=10pt,
  belowskip=10pt,
  language=Rust
}

% ---------------------------------------------------------------------------
% Custom environments (shared with the core specification)
% ---------------------------------------------------------------------------
\newtcolorbox{openquestion}[1][]{
  enhanced, breakable,
  colback=epiphanymist, colframe=epiphanycrimson,
  fonttitle=\bfseries\color{white}, title={\scshape\hspace{2pt}Open Question},
  coltitle=white, colbacktitle=epiphanycrimson,
  arc=1pt, boxrule=0pt, leftrule=2pt,
  left=10pt, right=10pt, top=8pt, bottom=8pt,
  attach boxed title to top left={xshift=0pt, yshift=0pt},
  boxed title style={arc=0pt, sharp corners, boxrule=0pt, left=6pt, right=8pt, top=2pt, bottom=2pt},
  #1
}
\newtcolorbox{rationale}[1][]{
  enhanced, breakable,
  colback=epiphanymist, colframe=epiphanyteal,
  fonttitle=\bfseries\color{white}, title={\scshape\hspace{2pt}Rationale},
  coltitle=white, colbacktitle=epiphanyteal,
  arc=1pt, boxrule=0pt, leftrule=2pt,
  left=10pt, right=10pt, top=8pt, bottom=8pt,
  attach boxed title to top left={xshift=0pt, yshift=0pt},
  boxed title style={arc=0pt, sharp corners, boxrule=0pt, left=6pt, right=8pt, top=2pt, bottom=2pt},
  #1
}
\newtcolorbox{requirement}[1][]{
  enhanced, breakable,
  colback=white, colframe=epiphanygold,
  fonttitle=\bfseries\color{white}, title={\scshape\hspace{2pt}Requirement},
  coltitle=white, colbacktitle=epiphanygold,
  arc=1pt, boxrule=0pt, leftrule=2pt,
  left=10pt, right=10pt, top=8pt, bottom=8pt,
  attach boxed title to top left={xshift=0pt, yshift=0pt},
  boxed title style={arc=0pt, sharp corners, boxrule=0pt, left=6pt, right=8pt, top=2pt, bottom=2pt},
  #1
}
\newtcolorbox{nongoal}[1][]{
  enhanced, breakable,
  colback=epiphanymist, colframe=epiphanyslate,
  fonttitle=\bfseries\color{white}, title={\scshape\hspace{2pt}Non-Goal},
  coltitle=white, colbacktitle=epiphanyslate,
  arc=1pt, boxrule=0pt, leftrule=2pt,
  left=10pt, right=10pt, top=8pt, bottom=8pt,
  attach boxed title to top left={xshift=0pt, yshift=0pt},
  boxed title style={arc=0pt, sharp corners, boxrule=0pt, left=6pt, right=8pt, top=2pt, bottom=2pt},
  #1
}

\newcommand{\MUST}{\textbf{MUST}}
\newcommand{\MUSTNOT}{\textbf{MUST}\nobreak\ \textbf{NOT}}
\newcommand{\SHOULD}{\textbf{SHOULD}}
\newcommand{\SHOULDNOT}{\textbf{SHOULD}\nobreak\ \textbf{NOT}}
\newcommand{\MAY}{\textbf{MAY}}

\setlist[itemize]{topsep=2pt, itemsep=3pt, parsep=0pt}
\setlist[enumerate]{topsep=2pt, itemsep=3pt, parsep=0pt}
\setlist[description]{topsep=2pt, itemsep=5pt, parsep=0pt}
\AtBeginDocument{\color{epiphanyink}}

% ---------------------------------------------------------------------------
% Document
% ---------------------------------------------------------------------------
\begin{document}

\begin{titlepage}
  \thispagestyle{empty}
  \centering
  \vspace*{2.2in}
  {\color{epiphanygold}\rule{3in}{0.8pt}}\\[18pt]
  {\titlefont\fontsize{34pt}{38pt}\selectfont\color{epiphanyteal}\bfseries Epiphany}\\[10pt]
  {\Large\scshape\color{epiphanyslate}Reference Suite}\\[6pt]
  {\large\itshape\color{epiphanyslate}A companion to the Core Specification}\\[14pt]
  {\color{epiphanygold}\rule{3in}{0.8pt}}\\[24pt]
  {\normalsize\color{epiphanyink}Version 0.1.0 --- Phase 3 (the initial entry set: six scores, required at Minimal, declared for Standard)}\\[4pt]
  {\small\color{epiphanyslate}Normative for solver conformance claims}
  \vfill
\end{titlepage}

\tableofcontents

% ===========================================================================
\chapter{About This Companion}
\label{ch:about}

The \emph{Reference Suite} is a companion to the Epiphany Core
Specification. It fulfils the core specification's delegation in its
\sectionsc{Companion Specifications} appendix (the
\texttt{sec:deferred:companions} section), which charters this document as
``the collection of test scores against which solver conformance is
established, with per-tier inclusion, per-tier metric thresholds, and any
fixed-expectation tests. Versioned with this specification.''

This release (v0.1.0) delivers:

\begin{itemize}
  \item the \textbf{suite entry model} --- how entries name their test
    scores, how each solve is configured, what passing an entry and
    passing the suite mean, and how fixed-expectation tests work
    (Chapter~\ref{ch:model});
  \item the \textbf{v0.1 entry set}: six scores, every one required for
    Minimal-tier conformance and the same six constituting the
    Standard-tier subset (Chapter~\ref{ch:entries});
  \item a non-normative note on the reference implementation's suite
    harness (Chapter~\ref{ch:harness}).
\end{itemize}

This document does \emph{not} cover:

\begin{itemize}
  \item the definition, normalization, and default thresholds of the
    quality metrics --- those are the \emph{Quality Metric Catalog}
    companion's, and this document consumes them;
  \item performance conformance (edit traces, frame budgets) --- the
    \emph{Performance Reference Suite} companion's, per the core
    specification's explicit boundary;
  \item the reference solving algorithm --- the non-normative
    \emph{Reference Algorithm} companion's.
\end{itemize}

\section{Relationship to the Core Specification}
\label{sec:about:relationship}

The suite's charter is the core specification's Chapter~9
(\sectionsc{The Constraint Solver Interface}, the \texttt{ch:solver}
chapter), \sectionsc{Conformance: The Reference Suite} section
(\texttt{sec:solver:conformance}). That section's suite-entry requirement
fixes what every entry consists of --- a test score in canonical
\texttt{.musc} form; per-tier inclusion; per-tier metric thresholds; and
optional fixed-expectation tests, used sparingly --- and fixes the pass
rule: ``A solver claiming a given tier \MUST{} pass every entry required
at that tier. Failure on any single entry is conformance failure at the
claimed tier.''

Three more core anchors bind this document:

\begin{itemize}
  \item The \sectionsc{Cross-Implementation Conformance} requirement (in
    the core's solver-determinism section) enumerates the four conditions
    a conforming solver meets across the suite; this document's
    per-entry evaluation rule (Section~\ref{sec:model:pass}) is those
    conditions applied entry-by-entry.
  \item The \sectionsc{Reference Algorithm} section makes
    fixed-expectation tests the \emph{only} place the suite may force a
    particular layout; everywhere else, any algorithm within thresholds
    conforms. Section~\ref{sec:model:fixed} inherits that discipline.
  \item The \emph{Quality Metric Catalog} companion defines every metric,
    its normalization, and the default per-tier thresholds this suite's
    entries reference. This document never restates a threshold; it
    names the catalog's defaults and records per-entry overrides (v0.1:
    none).
\end{itemize}

Where this document and a ratified core requirement disagree, \textbf{the
core requirement governs} and the discrepancy is a defect in this
document. One deliberate v0.1 deviation from the charter's letter ---
naming scores by deterministic builder rather than shipping
\texttt{.musc} bundles --- is called out as such, with its open question,
in Section~\ref{sec:model:referencing}.

\section{Conformance}
\label{sec:about:conformance}

The entry set, entry construction rules, solve configurations, and tier
inclusions in this document are \textbf{normative}. A solver conformance
claim at a tier is a claim about \emph{this} suite at \emph{this}
version:

\begin{itemize}
  \item claiming tier $T$ means passing every entry
    (Section~\ref{sec:model:pass}) that Chapter~\ref{ch:entries} requires
    at tier $T$ --- failure on any single entry is conformance failure at
    the claimed tier;
  \item per the core specification, suite versions are tied to
    specification versions, and a conforming implementation \MUST{}
    declare which suite version it passes
    (Section~\ref{sec:model:versioning}).
\end{itemize}

% ===========================================================================
\chapter{The Suite Entry Model}
\label{ch:model}

\section{Score Referencing}
\label{sec:model:referencing}

\begin{requirement}
\label{req:refsuite:referencing}
A v0.1 suite entry names its test score by \textbf{reference-implementation
builder and seed}: a deterministic constructor exported by the reference
implementation's test kit (the \texttt{epiphany-testkit} crate), together
with any seed argument, reproduces the score graph bit-for-bit. Two
referencing forms are used:

\begin{itemize}
  \item a \textbf{seeded builder}: a public function taking a \texttt{u64}
    seed (e.g.\ \texttt{fixtures::ten\_measure\_single\_staff}, or a
    generator re-exported through the test kit), with the entry pinning
    the exact seed value;
  \item a \textbf{corpus name}: the \texttt{name} string of an entry of
    the test kit's tagged corpus (\texttt{corpus()}), whose builder takes
    no arguments and is deterministic by construction.
\end{itemize}

The constructed score graph --- not any serialized artifact of it --- is
the entry's test score. Implementations under test \MUST{} evaluate the
entry against a score graph identical to the one the named builder
produces with the named seed.
\end{requirement}

\begin{rationale}
The core charter describes each entry as ``a test score in canonical
\texttt{.musc} form.'' v0.1 deliberately references builders instead of
shipping bundles, because today \emph{the builders are the canonical
definition}: they are versioned, reviewed, deterministic (seeded
\texttt{SplitMix64} identity minting), asserted invariant-clean, and
reproducible bit-for-bit by anyone building the reference crates --- while
the \texttt{.musc} byte format is still absorbing schema-major evolution,
so a shipped bundle would rot faster than the builder that made it. A
corpus caution: the test kit tags corpus fixtures with an
\emph{eligibility-taxonomy} tier (\texttt{Common} / \texttt{Edge} /
\texttt{Torture}); that taxonomy is unrelated to solver conformance tiers
(Minimal / Standard / Advanced) and carries no normative weight in this
suite.
\end{rationale}

\begin{openquestion}
\textbf{Builder references versus shipped \texttt{.musc} bundles.} The
charter's letter --- entries in canonical \texttt{.musc} form --- is not met
by v0.1: builder-plus-seed is implementation-anchored, which makes the
suite awkward for an independent implementation that does not link the
reference test kit (it must re-derive the score graphs from the builders'
sources). A future revision \MAY{} ship canonical \texttt{.musc} bundles
for every entry, exactly as the charter describes, once schema-major
bytes are stable enough that shipped bundles do not rot; at that point
builder references would remain as the bundles' provenance record. Until
then, the deviation is deliberate and this open question owns it.
\end{openquestion}

\section{Solve Configuration}
\label{sec:model:solve-config}

A suite entry is only reproducible if the solve it prescribes is fully
specified: the same score under a different page geometry casts off into
different systems, and several metric axes are defined over the declared
content bounds.

\begin{requirement}
\label{req:refsuite:solve-config}
Each suite entry declares the full solve configuration it is evaluated
under:

\begin{itemize}
  \item the \textbf{page geometry}: page size and margins, in staff
    spaces. In v0.1 every entry uses the reference implementation's
    documented default --- A4 portrait at an 8\,mm staff height: page
    $105 \times 148.5$ staff spaces, margins $7.5$ staff spaces on all
    four sides, hence a content area of $90 \times 133.5$ staff spaces.
    (The score graph has no home for page geometry yet: the core names
    \texttt{Canvas.layout\_defaults} without defining a type, tracked as
    Pass-12 row P12-I7, so geometry is declared per entry as a solver
    parameter.)
  \item the \textbf{solver configuration}: the \texttt{SolverConfig}
    fields. In v0.1 every entry uses the default configuration --- the
    \texttt{Standard} profile, an unbounded deterministic budget, and
    the Quality Metric Catalog's default tie-breaking weights.
\end{itemize}

An implementation \MUSTNOT{} substitute its own defaults for a declared
configuration when evaluating an entry.
\end{requirement}

\section{Passing an Entry, Passing the Suite}
\label{sec:model:pass}

\begin{requirement}
\label{req:refsuite:pass}
A solver \textbf{passes an entry at tier $T$} when, solving the entry's
score under the entry's declared configuration, all four of the core
specification's cross-implementation conformance conditions hold for that
solve:

\begin{enumerate}
  \item every hard constraint of the solve is satisfied;
  \item the result is internally deterministic per the core's
    within-implementation rule (byte-identical
    \texttt{ResolvedLayoutIR} for repeated identical solves within one
    implementation version);
  \item the \texttt{SolveReport} is well-formed and diagnostically
    accurate --- in particular the metric vector is computed per the
    Quality Metric Catalog, never a placeholder;
  \item every normative quality metric is at or below tier $T$'s
    threshold for its axis --- the Quality Metric Catalog's default
    per-tier threshold table, unless the entry declares a per-entry
    override (no v0.1 entry does);
\end{enumerate}

and additionally every fixed-expectation test the entry declares is
reproduced exactly (no v0.1 entry declares any).

A solver \textbf{passes the suite at tier $T$} when it passes every entry
required at tier $T$. Per the core specification, failure on any single
entry is conformance failure at the claimed tier; a status of
\texttt{Unsatisfiable} or a budget-exhausted partial solve on a suite
entry is a failure of condition~(1) or~(3), not an exemption.
\end{requirement}

\section{Fixed-Expectation Tests}
\label{sec:model:fixed}

The core specification permits an entry to pin specific layout properties
--- a particular bar's width, a specific system-break location --- that all
conforming solvers must reproduce, and directs that they be used
sparingly. In the resolved layout such expectations are checkable against
concrete structures: a bar's width against its \texttt{ResolvedMeasure}
bounding box, a system count or break location against the
\texttt{ResolvedSystem} list and the slots at which systems begin.

\begin{requirement}
\label{req:refsuite:fixed-expectations}
v0.1 declares \textbf{no fixed-expectation tests}. The mechanism is
deliberately unused: every v0.1 entry is evaluated by validity,
determinism, report accuracy, and thresholds alone. An entry of a future
suite revision that adds a fixed expectation \MUST{} state the expected
value, the resolved structure it is checked against, and the exact
comparison (including any tolerance).
\end{requirement}

\begin{rationale}
A fixed expectation binds \emph{every conforming solver} to one layout
fact forever after; it is the only place the suite may force a layout,
and none is warranted yet --- no cross-implementation ambiguity has
surfaced that thresholds fail to resolve. The reference implementation's
golden files (SVG snapshots, byte-anchored layouts) are
\emph{implementation regression locks}: they pin what \emph{that}
implementation produced so its own drift is caught, and they are
deliberately not suite conformance --- promoting them to fixed
expectations would freeze the ecosystem to the reference algorithm's
choices, exactly what the core's threshold-based conformance model
exists to avoid.
\end{rationale}

\section{Suite Versioning}
\label{sec:model:versioning}

\begin{requirement}
\label{req:refsuite:versioning}
Per the core specification, suite versions are tied to specification
versions: each release of this companion names the core-specification
version whose conformance story it serves, and a revision of the core's
solver chapter that changes tier obligations requires a corresponding
suite revision. This v0.1.0 suite serves the pre-1.0 core working draft.
A conforming implementation \MUST{} declare which suite version it
passes; a claim without a suite version is not a conformance claim.
\end{requirement}

This document keeps its own revision history (Chapter~\ref{ch:history}),
independent of the Quality Metric Catalog's: entry-set growth and
threshold-override changes are suite revisions, while metric definitions
and default thresholds revise in the catalog.

% ===========================================================================
\chapter{The v0.1 Entry Set}
\label{ch:entries}

\section{Overview}
\label{sec:entries:overview}

\begin{requirement}
\label{req:refsuite:entries}
The v0.1 entry set is exactly the six entries of
Table~\ref{tab:entries}. Every entry is \textbf{required at the Minimal
tier}. The \textbf{Standard-tier subset is the same six entries},
evaluated against the Quality Metric Catalog's Standard threshold
column. No entry declares a per-entry threshold override; no entry
declares a fixed-expectation test. No entry is designated
Advanced-only in v0.1 (the Advanced tier adds obligations, not
entries, until an extension's layout requirements enter the suite).
\end{requirement}

\begin{table}[h]
\centering
\small
\begin{tabular}{lp{4.2cm}p{6.6cm}}
\toprule
\textbf{Id} & \textbf{Entry} & \textbf{Construction} \\
\midrule
RS-1 & Ten-measure single staff &
  \texttt{fixtures::ten\_measure\_single\_staff} with seed
  \texttt{0x000A\_11CE} \\
RS-2 & Rich multi-region score &
  \texttt{generators::valid\_score\_rich} with seed \texttt{0xF302}
  (= corpus entry \texttt{gen\_valid\_score\_rich}) \\
RS-3 & B-flat major scale &
  corpus entry \texttt{b\_flat\_major\_scale} \\
RS-4 & Two-voice counterpoint &
  corpus entry \texttt{two\_voice\_counterpoint} \\
RS-5 & Notes and rests &
  corpus entry \texttt{notes\_and\_rests} \\
RS-6 & Three-four meter line &
  corpus entry \texttt{meter\_three\_four} \\
\bottomrule
\end{tabular}
\caption{The v0.1 entry set. All six entries are required at Minimal;
the same six constitute the Standard subset. Solve configuration for
every entry: the declared default geometry and solver configuration of
Requirement~\ref{req:refsuite:solve-config}.}
\label{tab:entries}
\end{table}

No implementation claims Standard-tier conformance as of this suite
version; the Standard listing exists so that the first Standard claim is
made against a pre-declared bar rather than a bar negotiated after the
fact. Where an entry's layout lacks a metric's geometry class (a single
system, a single page, no drawn slurs or beams), the affected axes
evaluate to $0.0$ under the catalog's vacuous-geometry rule and the
entry's threshold on those axes is trivially met; each entry's coverage
note below says which axes it exercises non-degenerately.

\section{RS-1: Ten-Measure Single Staff}
\label{sec:entries:rs1}

\begin{description}
  \item[Construction.]
    \texttt{epiphany\_testkit::fixtures::ten\_measure\_single\_staff(0x000A\_11CE)}
    --- the seed the reference implementation's engraving and rendering
    acceptance goldens document.
  \item[Content.] A 10-measure, single-staff, single-voice metric score:
    40 quarter notes (four per measure, all C4), plus a tie, a spanner, a
    marker, and a chord symbol. Invariant-clean by construction.
  \item[Coverage.] The multi-system workhorse. Under the declared
    geometry the spaced line (about 99 staff spaces) exceeds the content
    width of 90 staff spaces and casts off into \textbf{two systems},
    chosen among the barline break candidates --- so the
    \texttt{SystemBreak}, \texttt{CastingOff}, and
    \texttt{SymbolDensity} axes are exercised non-degenerately,
    alongside \texttt{Collision} and \texttt{Spacing} (axis names per
    the catalog's \texttt{QualityMetricKind} mapping). The
    cross-cutting objects (tie, spanner, marker, chord symbol) ride
    through the projection. Single page: the \texttt{PageFill} axis
    degenerates.
  \item[Tiers, thresholds, expectations.] Required at Minimal; in the
    Standard subset. Catalog default thresholds; no overrides; no fixed
    expectations.
\end{description}

\section{RS-2: Rich Multi-Region Score}
\label{sec:entries:rs2}

\begin{description}
  \item[Construction.]
    \texttt{generators::valid\_score\_rich(0xF302)} (the core crate's
    generator), identically reachable as the test-kit corpus entry
    \texttt{gen\_valid\_score\_rich} (the corpus pins the same seed).
  \item[Content.] Three \emph{concurrent} regions on disjoint staves: a
    metric region (measures, an eighth-note triplet tuplet, a tie, a
    spanner, a marker, a chord symbol, a decomposition attachment), a
    proportional region (wall-clock events), and an aleatoric region
    (musical-discipline events) --- plus tombstoned pitch and event ids
    and a spelling attachment resolving to a tombstoned pitch. Every
    core graph invariant holds.
  \item[Coverage.] Multi-region, multi-staff validity: concurrent
    regions must each be cast off and placed without cross-region
    interference. At the Minimal tier this entry demands validity and
    honest diagnostics on the projection it induces --- the non-metric
    regions reach the solver as ordinary constrained IR (their
    decompositions deferred upstream), so the entry does \emph{not}
    smuggle the Advanced tier's proportional/aleatoric layout obligation
    into Minimal; it guards that a multi-region document neither breaks
    hard constraints nor corrupts the report.
  \item[Tiers, thresholds, expectations.] Required at Minimal; in the
    Standard subset. Catalog default thresholds; no overrides; no fixed
    expectations.
\end{description}

\section{RS-3: B-flat Major Scale}
\label{sec:entries:rs3}

\begin{description}
  \item[Construction.] Corpus entry \texttt{b\_flat\_major\_scale}
    (zero-argument deterministic builder).
  \item[Content.] Seven quarter notes ascending through B-flat major
    (B$\flat$ C D E$\flat$ F G A) in one metric region, one staff, one
    voice --- accidentals inferred in a flat context.
  \item[Coverage.] The accidental entry: flats produce glyph clusters
    with left overhang, exercising overhang-aware spacing, the same-column
    cluster of accidental and notehead (the collision axis's same-slot
    exclusion), and the \texttt{Spacing} axis under uneven column
    ink. Single system and page: the break-family and page axes
    degenerate.
  \item[Tiers, thresholds, expectations.] Required at Minimal; in the
    Standard subset. Catalog default thresholds; no overrides; no fixed
    expectations.
\end{description}

\section{RS-4: Two-Voice Counterpoint}
\label{sec:entries:rs4}

\begin{description}
  \item[Construction.] Corpus builder \texttt{two\_voice\_counterpoint}
    (zero-argument, deterministic).
  \item[Content.] Two voices in one staff instance: an upper quarter-note
    line and a lower quarter-note line two octaves beneath it, sounding
    simultaneously.
  \item[Coverage.] Simultaneity: both voices' notes share musical time
    columns, so each spring slot carries a two-voice cluster, and the
    low line sits below the staff (ledger territory). Exercises
    same-column stacking, the collision sweep across a vertically spread
    texture, and column-advance regularity when columns are ink-heavy.
    Single system and page: the break-family and page axes degenerate.
  \item[Tiers, thresholds, expectations.] Required at Minimal; in the
    Standard subset. Catalog default thresholds; no overrides; no fixed
    expectations.
\end{description}

\section{RS-5: Notes and Rests}
\label{sec:entries:rs5}

\begin{description}
  \item[Construction.] Corpus entry \texttt{notes\_and\_rests}
    (zero-argument deterministic builder).
  \item[Content.] Note, rest, note, rest --- quarter values in one metric
    region, one staff, one voice.
  \item[Coverage.] Rest glyphs interleaved with noteheads: rest columns
    carry different ink boxes than note columns, exercising the spacing
    pass's treatment of mixed column content and the collision sweep over
    non-notehead glyphs. Single system and page: the break-family and
    page axes degenerate.
  \item[Tiers, thresholds, expectations.] Required at Minimal; in the
    Standard subset. Catalog default thresholds; no overrides; no fixed
    expectations.
\end{description}

\section{RS-6: Three-Four Meter Line}
\label{sec:entries:rs6}

\begin{description}
  \item[Construction.] Corpus entry \texttt{meter\_three\_four}
    (zero-argument deterministic builder).
  \item[Content.] Three quarter notes under a declared
    $\tfrac{3}{4}$ time signature (three quarter-beat groups), in one
    metric region, one staff, one voice.
  \item[Coverage.] Meter variety: the declared time signature drives
    measure length past the whole-note default, exercising the
    meter-resolution path that feeds the layout projection. Guards that
    a non-default meter neither breaks hard constraints nor perturbs
    spacing regularity. Single system and page: the break-family and
    page axes degenerate.
  \item[Tiers, thresholds, expectations.] Required at Minimal; in the
    Standard subset. Catalog default thresholds; no overrides; no fixed
    expectations.
\end{description}

% ===========================================================================
\chapter{Harness Binding (Non-Normative)}
\label{ch:harness}

This chapter is informative. The reference implementation binds this
suite to executable checks in its test kit, following the test kit's
established library-module-per-harness pattern (one public module per
harness, mirroring the corpus and prepass harnesses, with an integration
test driving it). The suite harness constructs each Chapter~\ref{ch:entries}
entry from its named builder and seed, solves it under the declared
configuration, and asserts the Minimal-tier pass of
Requirement~\ref{req:refsuite:pass}; it is delivered with the reference
implementation.

\begin{nongoal}
The harness is not part of this document's normative surface. Conformance
is defined by Chapters~\ref{ch:model} and~\ref{ch:entries} alone; an
independent implementation may bind the suite with any machinery that
evaluates the same entries under the same configurations. Likewise, the
reference implementation's golden files (SVG snapshots, byte-anchored
layouts) are that implementation's regression locks, not suite
conformance (Section~\ref{sec:model:fixed}).
\end{nongoal}

% ===========================================================================
\chapter{Revision History}
\label{ch:history}

\begin{longtable}{p{2cm} p{2.5cm} p{9cm}}
  \toprule
  \textbf{Date} & \textbf{Section} & \textbf{Change} \\
  \midrule
  \endhead
  \today & All & 0.1.0 --- Initial companion: pins the suite entry model
  (builder-and-seed score referencing with the \texttt{.musc}-bundle open
  question, per-entry solve configuration over the declared A4/8\,mm
  default geometry, the four-condition per-entry pass rule, the
  deliberately empty fixed-expectation set, suite-version declaration);
  delivers the six-entry v0.1 set
  (\texttt{ten\_measure\_single\_staff} seed \texttt{0x000A\_11CE},
  \texttt{valid\_score\_rich} seed \texttt{0xF302},
  \texttt{b\_flat\_major\_scale}, \texttt{two\_voice\_counterpoint},
  \texttt{notes\_and\_rests}, \texttt{meter\_three\_four}), all required
  at Minimal and all constituting the Standard subset under the Quality
  Metric Catalog's default thresholds, with no per-entry overrides. \\
\bottomrule
\end{longtable}

\end{document}
